\subsection{Measure-Preserving Dynamics and Ergodicity}

Measure-preserving dynamics study maps whose pullbacks keep mass fixed. The central question is whether time averages along one orbit recover space averages with respect to an invariant probability measure.

\begin{definition} \label{def:measure-preserving-and-ergodic-transformation} \hlt{Measure-Preserving and Ergodic Transformations}\\
Let $(X,\mathcal{M},\mu)$ be a measure space and $T:X\to X$.\\
For $E,A,B\in\mathcal{M}$, write $T^{-1}E\defeq\{x\in X:T(x)\in E\}$ and $A\triangle B\defeq(A\setminus B)\cup(B\setminus A)$.
\begin{enumerate}[label=(\roman*)]
\setlength{\itemsep}{0pt}
\item $T$ is \hlt{measurable} iff $T^{-1}E\in\mathcal{M}$ for all $E\in\mathcal{M}$.
\item A measurable $T$ is \hlt{measure preserving} iff $\mu(T^{-1}E)=\mu(E)$ for all $E\in\mathcal{M}$.
\item For measurable $T$, a set $A\in\mathcal{M}$ is \hlt{invariant modulo $\mu$} iff $\mu(A\triangle T^{-1}A)=0$.
\item If $\mu(X)=1$, $T$ measure preserving, then $T$ \hlt{ergodic} iff $\mu(A\triangle T^{-1}A)=0$ implies $\mu(A)\in\{0,1\}$ $\forall A\in\mathcal{M}$.
\end{enumerate}
\hlt{Key Insight}: Pullbacks preserve mass; ergodic invariants are trivial.
\end{definition}

\begin{proposition} \label{prop:measure-preserving-integral-test} \hlt{Integral Test for Measure Preservation}\\
Let $(X,\mathcal{M},\mu)$ be a finite measure space and let $T:X\to X$ be measurable.\\
Then $T$ is measure preserving iff $\int_X g\circ T\,d\mu=\int_X g\,d\mu$ for all $g\in L^1(X,\mu)$.
\begin{proof}
\begin{enumerate}[label=\roman*.]
\setlength{\itemsep}{0pt}
\item[$(\Rightarrow)$] Assume $T$ is measure preserving. For $A\in\mathcal{M}$, $\int_X\mathbf{1}_A\circ T\,d\mu=\mu(T^{-1}A)=\mu(A)=\int_X\mathbf{1}_A\,d\mu$. Thus $s=\sum_{k=1}^n c_k\mathbf{1}_{A_k}$ satisfies $\int_X s\circ T\,d\mu=\int_X s\,d\mu$ by linearity. If $0\le s_n\uparrow g$, then $0\le s_n\circ T\uparrow g\circ T$, so $\int_X g\circ T\,d\mu=\lim_n\int_X s_n\circ T\,d\mu=\lim_n\int_X s_n\,d\mu=\int_X g\,d\mu$. For $g\in L^1(X,\mu)$, apply this to $g^\pm$.
\item[$(\Leftarrow)$] For $A\in\mathcal{M}$, $\mu(X)<\infty$; $\mathbf{1}_A\in L^1(X,\mu)$, and $\mu(T^{-1}A)=\int_X\mathbf{1}_{T^{-1}A}\,d\mu=\int_X\mathbf{1}_A\circ T\,d\mu=\int_X\mathbf{1}_A\,d\mu=\mu(A)$.
\end{enumerate}
\hlt{Structural Commentary}: Indicators encode sets; approximation extends the identity to observables.\\
\hlt{Key Insight}: Measure preservation is equality of integrals after pullback.
\end{proof}
\end{proposition}

\begin{remark} \label{rem:invariant-set-indicator-test} \hlt{Invariant Sets as Invariant Indicators}\\
For a measure-preserving $T$, a measurable set $A$ is invariant modulo $\mu$ if and only if $\mathbf{1}_A\circ T=\mathbf{1}_A$ $\mu$-a.e.
\end{remark}

\begin{proposition} \label{prop:ergodicity-invariant-functions-constant} \hlt{Invariant Functions Detect Ergodicity}\\
Let $(X,\mathcal{M},\mu)$ be a probability space and let $T:X\to X$ be measure preserving. Then $T$ is ergodic iff every measurable $g:X\to\mathbb{R}$ with $g\circ T=g$ $\mu$-a.e. satisfies $g=c$ $\mu$-a.e. for some $c\in\mathbb{R}$.
\begin{proof}
\begin{enumerate}[label=\roman*.]
\setlength{\itemsep}{0pt}
\item[$(\Rightarrow)$] Assume $T$ is ergodic and $g\circ T=g$ $\mu$-a.e. For $a\in\mathbb{Q}$ set $A_a\defeq\{g<a\}$. Since $\mathbf{1}_{A_a}\circ T\ne\mathbf{1}_{A_a}$ only where $g\circ T\ne g$, $\mu(A_a\triangle T^{-1}A_a)=0$, so $\mu(A_a)\in\{0,1\}$. Since $A_n\uparrow X$ and $A_{-n}\downarrow\varnothing$, some $a_0,b_0\in\mathbb{Q}$ satisfy $\mu(A_{a_0})=0$ and $\mu(A_{b_0})=1$. Let $c\defeq\inf\{a\in\mathbb{Q}:\mu(A_a)=1\}$. For rational $a<c<b$, $\mu(A_a)=0$ and $\mu(X\setminus A_b)=0$; by countable unions, $\mu(g<c)=\mu(g>c)=0$, hence $g=c$ $\mu$-a.e.
\item[$(\Leftarrow)$] Assume $g\circ T=g$ $\mu$-a.e. implies $g$ is constant $\mu$-a.e. for every measurable $g:X\to\mathbb{R}$. If $A\in\mathcal{M}$ and $\mu(A\triangle T^{-1}A)=0$, then $\mathbf{1}_A\circ T=\mathbf{1}_{T^{-1}A}=\mathbf{1}_A$ $\mu$-a.e. Thus $\mathbf{1}_A$ is constant $\mu$-a.e., so $\mu(A)\in\{0,1\}$.
\end{enumerate}
\hlt{Structural Commentary}: Level sets convert invariant functions into invariant sets.\\
\hlt{Key Insight}: Ergodicity collapses invariant observables to constants.
\end{proof}
\end{proposition}

\begin{method} \label{meth:krylov-bogoliubov-averaging} \hlt{Krylov--Bogoliubov Averaging}\\
Let $X\ne\varnothing$ be compact metric, $f\in C(X,X)$, and $\mathcal{P}(X)\defeq\{\text{Borel probability measures on }X\}$.
\begin{enumerate}[label=(\roman*)]
\setlength{\itemsep}{0pt}
\item Choose $\eta\in\mathcal{P}(X)$ and set $\mu_n(E)\defeq\frac1n\sum_{k=0}^{n-1}\eta((f^k)^{-1}E)$ for $E\in\mathcal{B}(X)$.
\item Equivalently use $\int_X g\,d\mu_n=\int_X\frac1n\sum_{k=0}^{n-1}g\circ f^k\,d\eta$ for $g\in C(X)$.
\item Extract a weak-* convergent subsequence $\mu_{n_j}\to\mu$ in $\mathcal{P}(X)$.
\item Verify invariance by $\int_X g\circ f\,d\mu-\int_X g\,d\mu=\lim_j\frac1{n_j}\int_X(g\circ f^{n_j}-g)\,d\eta=0$ for $g\in C(X)$.
\end{enumerate}
For fixed $g\in C(X)$, $\abs{\frac1n\int_X(g\circ f^n-g)\,d\eta}\le2\|g\|_\infty/n$; weak-* limits are invariant-measure candidates.
\end{method}

\begin{figure}[H]
  \centering
  \begin{tikzpicture}[
  >=Stealth,
  every node/.style={font=\scriptsize},
  orbit/.style={draw=blue!55!black, fill=blue!5, rounded corners=1.5pt,
    minimum width=2.35cm, minimum height=0.68cm, align=center,
    inner xsep=6pt, inner ysep=4pt},
  measure/.style={draw=green!45!black, fill=green!8, rounded corners=1.5pt,
    minimum width=2.55cm, minimum height=0.68cm, align=center,
    inner xsep=6pt, inner ysep=4pt},
  compact/.style={draw=red!55!black, fill=red!5, rounded corners=1.5pt,
    minimum width=2.40cm, minimum height=0.68cm, align=center,
    inner xsep=6pt, inner ysep=4pt},
  arrow/.style={->, line width=0.55pt, black!65, shorten <=2pt, shorten >=2pt},
  dasharrow/.style={->, dashed, line width=0.5pt, black!55, shorten <=2pt, shorten >=2pt}
]
  \def\nodegap{0.86cm}
  \def\verticalgap{0.56cm}

  \node[orbit, anchor=west] (map) at (0,0) {$f:X\to X$\\continuous};
  \path (map.east) ++(\nodegap,0)
    node[measure, anchor=west] (avg) {orbit averages\\$\mu_n=\frac1n\sum_{k<n}(f^k)_*\eta$};
  \path (avg.east) ++(\nodegap,0)
    node[measure, anchor=west] (limit) {weak-* limit\\$\mu_{n_j}\to\mu$};
  \path (limit.east) ++(\nodegap,0)
    node[measure, anchor=west] (invariant) {invariant\\$f_*\mu=\mu$};

  \path (limit.south) ++(0,-\verticalgap)
    node[compact, anchor=north] (compactness) {compact metric $X$\\weak-* subsequence};

  \draw[arrow] (map.east) -- (avg.west);
  \draw[arrow] (avg.east) -- (limit.west);
  \draw[arrow] (limit.east) -- (invariant.west);

  \draw[dasharrow] (compactness.north) -- (limit.south);
\end{tikzpicture}

  \caption{Krylov--Bogoliubov: $\mu_n=\frac1n\sum_{k<n}(f^k)_*\eta\xrightarrow{w^*}\mu$ and $\frac1n(g\circ f^n-g)\to0$, hence $f_*\mu=\mu$.}
  \label{fig:krylov-bogoliubov-averaging-flow}
\end{figure}

\begin{theorem} \label{thm:bogoliubov-krilov-invariant-measure} \hlt{Bogoliubov--Krilov Invariant Measure Theorem}\\
If $X\ne\varnothing$ is compact metric and $f\in C(X,X)$, then $\exists\mu\in\mathcal{P}(X)$ with $f_*\mu=\mu$.
\begin{proof} Choose $\eta\in\mathcal{P}(X)$ and set $\mu_n\defeq\frac1n\sum\limits_{k=0}^{n-1}(f^k)_*\eta\in\mathcal{P}(X)$.\\
By Corollary \ref{cor:compact-metric-bounded-measures-weak-star-subsequence}, some $\mu_{n_j}\to\mu\in\mathcal{P}(X)$ weak-*; equivalently, $\int_X g\,d\mu_{n_j}\to\int_X g\,d\mu$ $\forall g\in C(X)$.\\
For $g\in C(X)$, $\int_X g\circ f\,d\mu_n-\int_X g\,d\mu_n=\frac1n\int_X(g\circ f^n-g)\,d\eta$, so $\abs{\int_X g\circ f\,d\mu_n-\int_X g\,d\mu_n}\le2\norm{g}_\infty/n$.\\
Passing to $n=n_j\to\infty$, $\int_X g\circ f\,d\mu=\int_X g\,d\mu$ $\forall g\in C(X)$. Since compact-metric Borel probabilities are Radon and $C_c(X)=C(X)$, by Theorem \ref{thm:radon-measures-determined-by-cc-integrals}, $f_*\mu=\mu$, i.e. $\mu(f^{-1}E)=\mu(E)$ $\forall E\in\mathcal{B}(X)$.\\
\hlt{Structural Commentary}: Cesaro pushforwards turn orbit motion into endpoint cancellation.\\
\hlt{Key Insight}: Weak-* limits kill the $1/n$ endpoint error.
\end{proof}
\end{theorem}

\begin{proposition} \label{prop:ergodic-measures-extreme-invariant-measures} \hlt{Ergodic Measures Are Extreme Invariant Measures}\\
Let $X$ be compact metric, $f\in C(X,X)$, $\mathcal{P}(X)$ Borel probabilities on $X$, $\mathcal{I}_f(X)\defeq\{\lambda\in\mathcal{P}(X):f_*\lambda=\lambda\}$.\\
For $\mu\in\mathcal{I}_f(X)$, the following are equivalent:
\begin{enumerate}[label=(\roman*)]
\setlength{\itemsep}{0pt}
\item $\mu\in\operatorname{Ext}(\mathcal{I}_f(X))$;
\item $f$ is $\mu$-ergodic.
\end{enumerate}
\begin{proof}
\begin{enumerate}[label=(\roman*)]
\setlength{\itemsep}{0pt}
\item[$(i)$]$\Rightarrow (ii)$ Assume not (ii): let $A\in\mathcal{B}(X)$, $\mu(A\triangle f^{-1}A)=0$, $0<\mu(A)<1$; $\nu(E)\defeq\frac{\mu(E\cap A)}{\mu(A)}$, $\rho(E)\defeq\frac{\mu(E\cap A^c)}{\mu(A^c)}$.\\
For $E\in\mathcal{B}(X)$, $\mu(f^{-1}E\cap A)=\mu(f^{-1}E\cap f^{-1}A)=\mu(f^{-1}(E\cap A))=\mu(E\cap A)$, similarly for $A^c$.\\
Hence $\nu,\rho\in\mathcal{I}_f(X)$. Since $\mu=\mu(A)\nu+(1-\mu(A))\rho$, $\nu(A)=1$, $\rho(A)=0$; $\mu$ is not extreme.
\item[$(ii)$]$\Rightarrow (i)$ Let $\mu=\theta\nu+(1-\theta)\rho$ with $0<\theta<1$ and $\nu,\rho\in\mathcal{I}_f(X)$.\\
Since $\theta\nu\leq\mu$, by Theorem \ref{thm:radon-nikodym-sigma-finite}, $h=\frac{d\nu}{d\mu}$ with $0\leq h\leq\frac1\theta$ $\mu$-a.e.\\
Then $\int_X h\circ f\,h\,d\mu=\int_X h\circ f\,d\nu=\int_X h\,d\nu=\int_X h^2\,d\mu$, while $\int_X(h\circ f)^2\,d\mu=\int_Xh^2\,d\mu$ by $\mu$-invariance.\\
Thus $\int_X(h\circ f-h)^2\,d\mu=0$, so $h\circ f=h$ $\mu$-a.e. By Proposition \ref{prop:ergodicity-invariant-functions-constant}, $h=c$ $\mu$-a.e.\\
Since $1=\nu(X)=\int_Xh\,d\mu=c$, $\nu=\mu$, and then $\rho=\mu$.
\end{enumerate}
\hlt{Structural Commentary}: Invariant sets split measures; invariant densities test splits.\\
\hlt{Key Insight}: Extreme invariant measures are exactly ergodic measures.
\end{proof}
\end{proposition}

\begin{corollary} \label{cor:compact-map-ergodic-invariant-measure} \hlt{Existence of an Ergodic Invariant Measure}\\
Let $X\ne\varnothing$ be compact metric, $f\in C(X,X)$, and $\mathcal{I}_f(X)\defeq\{\mu\in\mathcal{P}(X):f_*\mu=\mu\}$.\\
Then $\exists\mu\in\mathcal{I}_f(X)$ such that $f$ is $\mu$-ergodic.
\begin{proof} By Theorem \ref{thm:bogoliubov-krilov-invariant-measure}, $\mathcal{I}_f(X)\ne\varnothing$. Convexity immediate. If $\mu_\alpha\to\mu$ weak-*, $\mu_\alpha\in\mathcal{I}_f(X)$, then $\mu\in\mathcal{P}(X)$ and $\int_X g\circ f\,d\mu=\lim_\alpha\int_X g\circ f\,d\mu_\alpha=\lim_\alpha\int_X g\,d\mu_\alpha=\int_X g\,d\mu$ $\forall g\in C(X)$.\\
By Theorem \ref{thm:radon-measures-determined-by-cc-integrals}, $f_*\mu=\mu$, so $\mu\in\mathcal{I}_f(X)$. Thus $\mathcal{I}_f(X)$ is weak-* closed and $\norm{\mu}_{\mathrm{var}}=1$ on $\mathcal{I}_f(X)$.\\
By Corollary \ref{cor:weak-star-compact-radon-measure-sets}, $\mathcal{I}_f(X)$ is weak-* compact.\\
By Krein--Milman, $\exists$ extreme $\mu\in\mathcal{I}_f(X)$. By Proposition \ref{prop:ergodic-measures-extreme-invariant-measures}, $\mu$ is ergodic.\\
\hlt{Structural Commentary}: Compactness supplies an extreme invariant measure.\\
\hlt{Key Insight}: Extreme invariant measures are ergodic.
\end{proof}
\end{corollary}

\begin{theorem} \label{thm:ergodic-averages-characterization} \hlt{Asymptotic Ergodic Averages}\\
Let $T:X\to X$ be measure preserving on probability space $(X,\mathcal{M},\mu)$, $U_Tg\defeq g\circ T$, $A_n\defeq\frac1n\sum\limits_{k=0}^{n-1}U_T^k$.\\
Then $T$ is ergodic iff $A_ng\to\int_Xg\,d\mu$ $\mu$-a.e. $\forall g\in L^1(X,\mu)$.
If $0<\mu(X)<+\infty$, replace $\int_Xg\,d\mu$ by $\frac{1}{\mu(X)}\int_Xg\,d\mu$.
\begin{proof}
\begin{enumerate}[label=\roman*.]
\setlength{\itemsep}{0pt}
\item[$(\Rightarrow)$] By Birkhoff's pointwise ergodic theorem, $A_ng\to h$ $\mu$-a.e. for some $h\in L^1(X,\mu)$ with $h\circ T=h$ $\mu$-a.e. and $\int_Xh\,d\mu=\int_Xg\,d\mu$. Since $T$ is ergodic, Proposition \ref{prop:ergodicity-invariant-functions-constant} implies $h=c$ $\mu$-a.e.; by integrating, $c=\int_Xg\,d\mu$.
\item[$(\Leftarrow)$] Let $A\in\mathcal{M}$ with $\mu(A\triangle T^{-1}A)=0$. For $g=\mathbf{1}_A$, $A_n\mathbf{1}_A=\mathbf{1}_A$ $\mu$-a.e.\\
By the assumed limit, $\mathbf{1}_A=\mu(A)$ $\mu$-a.e., hence $\mu(A)\in\{0,1\}$.
\end{enumerate}
\hlt{Structural Commentary}: The ergodic limit is the invariant component; ergodicity makes it constant.\\
\hlt{Key Insight}: Time averages equal the space mean exactly when invariants are trivial.
\end{proof}
\end{theorem}

\begin{proposition} \label{prop:koopman-isometry-and-linear-measure-preservation} \hlt{Koopman and Linear Examples}
\begin{enumerate}[label=(\roman*)]
\setlength{\itemsep}{0pt}
\item If $\mu(X)<+\infty$, $T:X\to X$ is measurable, and $U_Tg\defeq g\circ T$,\\
then $T$ preserves $\mu$ iff $\|U_Tg\|_p=\|g\|_p$ $\forall g\in L^p(X,\mu)$ and $1\leq p<+\infty$.
\item For linear $T:\R^n\to\R^n$, $T_*m_n=m_n$ iff $T\in GL_n(\R)$ and $\abs{\det T}=1$; hence $\operatorname{diag}(2,\frac12)_*m_2=m_2$.
\end{enumerate}
\begin{proof}
\begin{enumerate}[label=(\roman*)]
\setlength{\itemsep}{0pt}
\item[$(i)$] ($\Rightarrow$) If $T$ preserves $\mu$, then by Proposition \ref{prop:measure-preserving-integral-test}, $\|U_Tg\|_p^p=\int_X\abs{g\circ T}^p\,d\mu=\int_X\abs{g}^p\,d\mu=\|g\|_p^p$.
\item[$(i)$] ($\Leftarrow$) If the $L^1$ isometry holds, then for $A\in\mathcal{M}$, $\mu(T^{-1}A)=\|\mathbf{1}_A\circ T\|_1=\|\mathbf{1}_A\|_1=\mu(A)$.
\item[$(ii)$] ($\Rightarrow$) If $T_*m_n=m_n$ and $T$ is singular, then $m_n(T(\R^n))=0$ by Corollary \ref{cor:linear-images-rigid-motion-lebesgue-measure}, but $T^{-1}(T(\R^n))=\R^n$, contradiction. Thus $T\in GL_n(\R)$. By Corollary \ref{cor:invertible-linear-images-lebesgue-sets}, $m_n(T^{-1}E)=\abs{\det T}^{-1}m_n(E)$.\\
Any $E$ with $0<m_n(E)<+\infty$ implies $\abs{\det T}=1$.
\item[$(ii)$] ($\Leftarrow$) If $T\in GL_n(\R)$ and $\abs{\det T}=1$, then by Corollary \ref{cor:invertible-linear-images-lebesgue-sets},\\
$m_n(T^{-1}E)=\abs{\det T}^{-1}m_n(E)=m_n(E)$ for all Lebesgue measurable $E$.
\end{enumerate}
\hlt{Structural Commentary}: Pullbacks preserve norms or scale by determinant.\\
\hlt{Key Insight}: Measure preservation is norm preservation or unit determinant.
\end{proof}
\end{proposition}

\begin{theorem} \label{thm:poincare-recurrence-measure-preserving} \hlt{Poincare Recurrence}\\
Let $\mu(X)<+\infty$, $T:X\to X$ preserve $\mu$, $A\in\mathcal{M}$.\\
With $R_A\defeq\{x\in A:\#\{n\geq1:T^nx\in A\}=+\infty\}$, $\mu(A\setminus R_A)=0$.
\begin{proof}
Set $E\defeq A\cap\bigcap\limits_{n\geq1}T^{-n}(A^c)$. Then $E\cap T^{-m}A=\varnothing$ for $m\geq1$.\\
If $x\in T^{-i}E\cap T^{-j}E$ with $0\leq i<j$, then $T^ix\in E\cap T^{-(j-i)}E\subseteq E\cap T^{-(j-i)}A$, contradiction; hence $(T^{-k}E)_{k\geq0}$ are disjoint. For $N\geq0$, by preservation and finite additivity,\\
$(N+1)\mu(E)=\mu(\bigcup\limits_{k=0}^NT^{-k}E)\leq\mu(X)$, so $\mu(E)\leq\frac{\mu(X)}{N+1}$ $\forall N$; hence $\mu(E)=0$.\\
If $x\in A\setminus R_A$ and $r\geq0$ is its last visit time, then $T^rx\in E$, so $A\setminus R_A\subseteq\bigcup\limits_{r\geq0}T^{-r}E$ and $\mu(A\setminus R_A)=0$.\\
\hlt{Structural Commentary}: The nonreturn set is wandering; finite-return points eventually enter it.\\
\hlt{Key Insight}: Finite measure forbids positive-measure wandering sets.
\end{proof}
\end{theorem}
